\chapter[REVISÃO DO ESTADO DA ARTE]{REVISÃO DO ESTADO DA ARTE}

A revisão de literatura deste capítulo está organizada em três eixos temáticos diretamente relacionados à proposta desta dissertação: (i) classificação automatizada de lesões cutâneas por aprendizado profundo, com ênfase na detecção de melanoma; (ii) aplicações de modelos de linguagem multimodal em diagnóstico dermatológico; e (iii) sistemas multiagentes para suporte à decisão clínica. Ao final, apresenta-se o posicionamento do trabalho proposto em relação à literatura revisada.


\section{Classificação Automatizada de Lesões Cutâneas}

A aplicação de redes neurais convolucionais ao diagnóstico de lesões cutâneas ganhou relevância expressiva com o estudo de \citeonline{esteva2017dermatologist}, que treinaram uma CNN de arquitetura Inception v3 em um conjunto de 129.450 imagens clínicas -- incluindo 3.374 dermatoscópicas --, das quais 127.463 foram usadas para treinamento e validação e 1.942 para teste, com diagnóstico confirmado por biópsia. A comparação foi realizada com pelo menos 21 dermatologistas certificados, vinculados aos departamentos de dermatologia de Stanford University, University of Pennsylvania, Massachusetts General Hospital e University of Iowa (25, 22 e 21 participantes nas comparações de carcinoma, melanoma clínico e melanoma dermatoscópico, respectivamente), em duas tarefas binárias -- carcinomas queratinocíticos versus ceratoses seborreicas benignas, e melanoma maligno versus nevo benigno --, contemplando tanto imagens clínicas fotográficas quanto dermatoscópicas. A CNN atingiu AUC de 0,96, 0,94 e 0,91 nas três comparações, desempenho equivalente ao dos dermatologistas (\textit{on par with all tested experts}). Esse resultado estabeleceu um marco na área e motivou investigações subsequentes sobre a capacidade de modelos de aprendizado profundo em tarefas dermatoscópicas.

\citeonline{haenssle2018man} aprofundaram essa análise ao comparar uma CNN de arquitetura Inception v4 com 58 dermatologistas de 17 países, convidados por meio de listas de e-mail da International Dermoscopy Society, dos quais 30 foram classificados como \textit{experts} (mais de cinco anos de experiência em dermatoscopia), na classificação de nevos melanocíticos e melanomas. O conjunto de teste reuniu 100 imagens dermatoscópicas, selecionadas por dois dermatologistas experientes entre 300 imagens por apresentarem maior dificuldade diagnóstica, e foi avaliado em dois níveis de informação: apenas dermatoscopia (nível I) e dermatoscopia associada a dados clínicos (nível II). Nessas condições, a CNN atingiu especificidade de 82,5\%, superior à dos dermatologistas nos níveis I e II (71,3\% e 75,7\%, respectivamente, para sensibilidades equivalentes de 86,6\% e 88,9\%), e AUC-ROC de 0,86, superior à área ROC média dos dermatologistas no nível I (0,79) e no nível II (0,82). Esses achados reforçam o potencial clínico dos modelos supervisionados, especialmente quando treinados em conjuntos de dados de grande escala e alta qualidade.

No domínio específico do melanoma, \citeonline{brinker2019deep} treinaram uma CNN com 12.378 imagens dermatoscópicas de código aberto e compararam seu desempenho ao de 157 dermatologistas de 12 hospitais universitários alemães, de diferentes níveis de experiência (de médicos em formação a médicos-chefes), em uma tarefa de classificação dermatoscópica de melanoma com um conjunto de teste de 100 imagens. Os dermatologistas apresentaram sensibilidade média de 74,1\% e especificidade média de 60\%; na especificidade média dos dermatologistas, a CNN atingiu sensibilidade média de 87,5\%, e, na sensibilidade média deles, especificidade média de 86,5\%. A CNN, treinada exclusivamente com imagens de código aberto, superou 136 dos 157 dermatologistas avaliados.

A disponibilização de conjuntos de dados públicos ampliou significativamente o ritmo de desenvolvimento nessa área. O conjunto HAM10000, introduzido por \citeonline{tschandl2018ham10000}, disponibilizou mais de 10 mil imagens dermatoscópicas anotadas abrangendo sete categorias de lesões cutâneas, tornando-se um dos principais \textit{benchmarks} da comunidade. Em paralelo, os desafios anuais promovidos pela \textit{International Skin Imaging Collaboration} (ISIC) consolidaram protocolos padronizados de avaliação e incentivaram o desenvolvimento de soluções baseadas em \textit{transfer learning} e \textit{ensemble} de modelos \cite{codella2019skin}. Nesse mesmo sentido, destaca-se a relevância de conjuntos com aquisição padronizada, como o PH2, que oferece 200 imagens dermatoscópicas com anotações especializadas e máscaras de segmentação de referência fornecidas por especialistas do Hospital Pedro Hispano, Portugal \cite{mendonca2013ph2}. A disponibilidade dessas máscaras torna o PH2 especialmente adequado para estudos que buscam isolar o componente diagnóstico da variabilidade introduzida por modelos de segmentação automática, como é o caso desta pesquisa.


\section{Modelos de Linguagem Multimodal em Diagnóstico Médico}

A evolução dos modelos de linguagem de larga escala para arquiteturas multimodais abriu novas perspectivas para o diagnóstico médico auxiliado por inteligência artificial. \citeonline{nori2023capabilitiesgpt4medicalchallenge} avaliaram as capacidades do GPT-4 em exames de certificação médica e tarefas clínicas estruturadas, demonstrando que modelos de propósito geral podem atingir desempenho relevante em raciocínio médico sem ajuste fino específico. No campo da visão computacional, \citeonline{liu2023visual} introduziram o LLaVA, um modelo multimodal baseado em instrução que combina \textit{encoders} visuais com modelos de linguagem, possibilitando a análise integrada de imagens e texto. Essa capacidade de generalização em contextos clínicos, somada à arquitetura de integração visual-textual do LLaVA, que permite descrever achados visuais em linguagem natural, motivou investigações mais específicas em diagnóstico por imagem e em domínios como a dermatologia.

Especificamente em dermatologia, \citeonline{vashisht2024umassbionlpmediqam3g2024dermprompt} avaliaram sistematicamente estratégias de engenharia de \textit{prompts} com GPT-4V para o diagnóstico dermatológico por meio do DermPrompt, demonstrando que a estruturação do raciocínio clínico no \textit{prompt} melhora a precisão diagnóstica e gera respostas mais auditáveis em relação a classificadores tradicionais de ponta a ponta. \citeonline{wang2023chatcadinteractivecomputeraideddiagnosis} propuseram o ChatCAD, um sistema de diagnóstico assistido por computador que incorpora LLMs para geração de relatórios em linguagem natural a partir de imagens médicas, melhorando a comunicação entre o sistema e os profissionais de saúde.


\section{Sistemas Multiagentes para Suporte Diagnóstico}

A aplicação de sistemas multiagentes ao diagnóstico médico representa uma abordagem natural para simular a dinâmica de consultoria clínica interdisciplinar, na qual múltiplos especialistas contribuem com perspectivas complementares para uma decisão final. \citeonline{ke2024mitigating} investigaram estratégias de mitigação de vieses em sistemas multiagentes aplicados à medicina, demonstrando que a diversificação de agentes e a incorporação de mecanismos críticos reduzem a propagação de erros individuais. A viabilidade dos SMAs baseados em LLMs em contextos clínicos tem sido demonstrada em trabalhos recentes. \citeonline{tang2024medagents} propuseram o MedAgents, no qual múltiplos agentes especializados em diferentes áreas médicas colaboram por meio de discussão estruturada para responder a questões de raciocínio clínico em regime \textit{zero-shot}, superando abordagens de agente único em benchmarks médicos. \citeonline{kim2024mdagents} avançaram essa direção com o MDAgents, um framework que aloca dinamicamente a estrutura de colaboração entre agentes (individual, par ou grupo) de acordo com a complexidade do caso clínico, obtendo ganhos de até 4,2\% em benchmarks de tomada de decisão médica frente a abordagens de agente único. Esses resultados indicam que a especialização de agentes para tarefas diagnósticas distintas, triagem, análise morfológica e geração de relatórios, proporciona maior robustez e interpretabilidade em relação a sistemas monolíticos, fundamentando teoricamente essa escolha arquitetural \cite{wooldridge2009introduction}. Essa estrutura favorece a integração de diferentes fontes de conhecimento e a melhoria da auditabilidade das decisões em sistemas de apoio clínico, características essenciais para aplicações em contextos de saúde de alto risco.


\section{Posicionamento deste Trabalho}

A literatura revisada evidencia um campo em rápida evolução, no qual modelos supervisionados de alta performance coexistem com abordagens emergentes baseadas em LLMs e arquiteturas multiagentes. Enquanto CNNs e modelos de \textit{transfer learning} estabeleceram resultados sólidos em \textit{benchmarks} como HAM10000 e ISIC, sua natureza de ``caixa-preta'' limita a interpretabilidade clínica das decisões. Por outro lado, LLMs multimodais oferecem diagnósticos explicáveis em linguagem natural, mas ainda apresentam desempenho quantitativo inferior aos classificadores especializados em tarefas binárias de detecção de melanoma.

Em dermatologia, a aplicação de arquiteturas multiagentes permite simular uma junta médica virtual, na qual agentes especializados em protocolos clínicos distintos oferecem pareceres independentes sobre a mesma lesão. A síntese dessas análises por um agente consolidador, complementada pela revisão crítica de um agente baseado em recuperação de conhecimento (\textit{Retrieval-Augmented Generation}, RAG), proporciona diagnósticos mais robustos e auditáveis, reduzindo a dependência de uma única perspectiva analítica.

Nesse cenário, trabalhos como MedAgents \cite{tang2024medagents} e MDAgents \cite{kim2024mdagents} demonstram o potencial de arquiteturas multiagentes baseadas em LLMs para o raciocínio clínico geral, mas operam sobre entradas textuais (questões e casos descritos em linguagem natural), sem lidar com o componente visual do diagnóstico nem com protocolos clínicos estruturados e ponderados como os empregados em dermatoscopia. Esta dissertação contribui para preencher essa lacuna ao avaliar, de forma sistemática, a capacidade diagnóstica multimodal de GPT-5.4-nano e Gemini-3.1-flash-lite em um cenário controlado: imagens dermoscópicas do conjunto PH2 com máscaras de segmentação de referência como entrada, os protocolos clínicos ABCD, Menzies e SPCL como estruturas de raciocínio diagnóstico, e arquitetura multiagente como mecanismo de consolidação de múltiplas perspectivas especializadas. Essa configuração isola o componente de raciocínio das LLMs das demais fontes de variabilidade do sistema, oferecendo uma avaliação interna válida da viabilidade dessa abordagem para a detecção de melanoma com diagnósticos explicáveis.

Os modelos avaliados nesta dissertação, GPT-5.4-nano \cite{openai2025gpt5nano} e Gemini-3.1-flash-lite \cite{google2025gemini31}, representam gerações mais recentes de modelos multimodais de larga escala, com capacidades de raciocínio e compreensão visual expandidas em relação a versões anteriores. A avaliação sistemática dessas arquiteturas em tarefas estruturadas de detecção de melanoma, utilizando os protocolos ABCD, Menzies e SPCL como estruturas de raciocínio clínico, constitui uma contribuição original deste trabalho.
